\documentclass[margin,line]{resume}
 
\usepackage[english, portuguese]{babel}
\usepackage[T1]{fontenc}
\usepackage{fontawesome5}
\usepackage{graphicx,wrapfig}
\usepackage{url}
\usepackage[colorlinks=false, pdfstartview=FitV, citecolor=black, urlcolor=black]{hyperref}
\usepackage{cite}
% tighter bibliography: no extra space between references
\makeatletter
\renewcommand{\@openbib@code}{\setlength{\itemsep}{1pt}\setlength{\parsep}{0pt}}
\makeatother
%\usepackage[backend=biber]{biblatex}
%\addbibresource{publications.bib}


\begin{document}{\large Taylla Milena Theodoro}
\begin{resume}


% === PERSONAL INFO ===

    \section{\mysidestyle Contact \\Information}
    \begin{minipage}[t]{0.7\textwidth}
        Computational Imaging \\
        Biohub SF\\
    \end{minipage}\hfill
    \begin{minipage}[t]{0.7\textwidth}
        \faIcon{envelope}
        \href{mailto: taylla.theodoro@czbiohub.org}
        {\underline{taylla.theodoro@czbiohub.org}} \\
        \faIcon{linkedin}
        \href{https://www.linkedin.com/in/taylla-theodoro/}{\underline{taylla-theodoro}} 
    \end{minipage}
   \vspace{-0.2cm}
        \section{\mysidestyle Research and Work Experience}
        \textbf{\href{https://biohub.org/comp-imaging/people/}{Computational Imaging - Chan Zuckerberg Biohub SF}} \hfill \textsl{Aug 2024 -- present}  \\
        \textit{R\&D Engineer} - Microscope automation and high-throughput analysis of dynamic live-cell imaging.
        \begin{list2}
            \item Built DynaTrack, a label-agnostic smart-tracking loop that runs phase reconstruction and virtual staining during acquisition to follow zebrafish neuromasts with gentle bright-field imaging; kept 73 of 75 neuromasts in view for up to 70.5\,h, versus $\sim$12\,h with classical imaging~\cite{theodoro2026dynatrack}.
            \item Curated the shared 5D (3D+t, multi-channel) live-cell dataset behind DynaCLR and DynaCell: registration, virtual-staining-based nuclear segmentation, joint segmentation and tracking with Ultrack, and OME-Zarr organization of infected and drug-perturbed cells. This dataset underpins DynaCLR's infection-state classifier (F1 98.4)~\cite{pradeep2024contrastivelearningcellstate,hirata2025mapping,hirata-miyasaki2025embedding}.
            \item Curated and registered the DynaCell benchmark for dynamic 3D virtual staining: paired label-free and fluorescence volumes of A549 cells across four organelles and three infection conditions, acquired on the Mantis microscope~\cite{kalinin2026dynacell}.
        \end{list2}
        \vspace{-0.2cm}

        \textbf{\href{https://lids.ic.unicamp.br/people}{LIDS - Laboratory of Image Data Science, UNICAMP}} \hfill \textsl{Aug 2021 -- 2025}  \\
        \textit{Graduate Researcher} - Automated detection of Early Contrast Enhancement (ECE), an MRI biomarker of pleural mesothelioma, in 4D (3D+t) DCE-MRI of the thorax.
        \begin{list2}
            \item Built Meso-Regions, an open-source (GPL-3) five-stage pipeline: ANTs motion correction, image-foresting-transform pleural extraction, SLIC/DISF supervoxel partitioning, per-supervoxel enhancement curves, and rule-based ECE detection. Flags suspicious regions with no manual ROI placement; published at IEEE ISBI 2026~\cite{11515662,theodoro2025mesothelioma}.
            \item Trained a PyTorch U-Net with diffusion-based data augmentation to segment the pleural region containing the tumor; co-authored a superpixel tutorial at SIBGRAPI and the iMig imaging review~\cite{9643120,ARMATO2024107832}.
        \end{list2}
        \vspace{-0.2cm}

        \textbf{\href{https://www.predictmeso.com/members/}{PREDICT-Meso Network, University of Glasgow}} \hfill \textsl{Feb 2023 -- Jul 2023}\\
        \textit{Visiting Researcher} - Built MISe, an interactive segmentation tool for 4D DCE-MRI, and used it to produce ground-truth pleural labels for the PREDICT-Meso cohort~\cite{mise2023}.
        \vspace{-0.3cm}

        \textbf{Peapply \& Co.} \hfill \textsl{May 2021 -- Oct 2021}\\
        \textit{Data Engineer} - Designed and built pipelines for collecting, storing, and analyzing data at scale in PySpark and SQL.
        \vspace{-0.3cm}

        \textbf{\href{https://sites.google.com/view/labiii/members/current?authuser=0}{LIII - Lab of Image, Intelligence and Innovation}} \hfill \textsl{Jan 2020 -- Apr 2020}\\
        \textit{Visiting Student Researcher} - CNN-based MRI spine segmentation in TensorFlow.
        \vspace{-0.3cm}

        \textbf{\href{https://www.linkedin.com/company/la2i---laborat-rio-de-automa-o-e-instrumenta-o-inteligente/}{LAB2I - Automation and Intelligent Instrumentation Laboratory}}  \hfill \textsl{Aug 2017 -- Dec 2019}\\
        \textit{Undergraduate Research Intern} - Built an embedded sensing system to monitor temperature and electric current in electrocoagulation water treatment~\cite{taylla2018,taylla2019}.
    \vspace{-0.2cm}
       \section{\mysidestyle Education}
        \textbf{University of Campinas (UNICAMP)} SP, BR  \hfill \textsl{Aug 2021 -- 2025}\\
           \textit{M.Sc. in Computer Science}, GPA 4.0/4.0
           \begin{list2}
            \item Thesis: Towards Automatic Detection of Pleural Mesothelioma Biomarker in 4D Dynamic MR Imaging~\cite{theodoro2025mesothelioma}
           \end{list2}
        \textbf{State University of Londrina (UEL)} PR, BR  \hfill \textsl{Apr 2016 -- Mar 2021}\\
           \textit{B.Sc. in Electrical Engineering}
           \begin{list2}
            \item Thesis: Study of Deep Learning and its Application to Medical Images
           \end{list2}

     \vspace{-0.3cm}
    \section{\mysidestyle Advanced Training}
    
    \textbf{Marine Biological Laboratory (MBL), Woods Hole, MA, USA} \hfill \textsl{2025}\\
    \textit{Analytical and Quantitative Light Microscopy}\\
    Advanced intensive research training course focused on quantitative fluorescence microscopy, optical physics, image analysis, and computational methods for biological imaging.
    \section{\mysidestyle Coursework}  
      \textbf{Graduate} Digital Image Processing, Image Analysis, Parallelization and Optimization Algorithms for Machine Learning Graphs, Machine Learning and Pattern Recognition,  Deep Learning, Graph Algorithms.
    \vspace{-0.3cm}

      \textbf{Undergraduate} Digital Circuits, Analog and Digital Electronics, C programming, Embedded Systems, Digital Signal Processing, Modern Programming Languages, Communication Principles, Electronic Instrumentation, Microelectronics (ARM), Biomedical Engineering, VHDL.
       \vspace{-0.3cm}

     \section{\mysidestyle Teaching Assistant}
             MC202 - Data Structure \hfill \textsl{Aug 2022 - Nov 2022}\\
             MC832 - Computer Network \hfill \textsl{Mar 2022 - Jun 2022}\\    
       \vspace{-0.3cm} 
     \section{\mysidestyle Volunteer Experience}
        Techninas - Girls in Tech \hfill \textsl{2018 - 2020}\\
        Student Branch IEEE UEL \hfill \textsl{2016 - 2018}\\
        Women in Engineering IEEE UEL \hfill \textsl{2016 - 2018}\\
    \vspace{-0.3cm}
      \section{\mysidestyle Conference Presentation}
   
        \textbf{ISBI 2026 - IEEE International Symposium on Biomedical Imaging} Meso-Regions: 4D Segmentation of Early Contrast Enhancement Biomarkers for Pleural Mesothelioma Diagnosis\cite{11515662}.\\ 
        \textbf{SPIE BiOS 2026} DynaTrack: label-agnostic smart tracking for mapping dynamics of organ development \cite{theodoro2026dynatrack}.\\
        \textbf{iMig 2023 - 16th International Conference of the International Mesothelioma interest group} Development of an automated machine learning tool for reporting of MRI-Early Contrast Enhancement in Mesothelioma\cite{hirata2025mapping}. \\
        \textbf{SIBGRAPI 2023 - 36th Conference on Graphics, Patterns and Images} Workshop - Superpixel segmentation: from theory to applications\cite{9643120}.\\
    %        \section{\mysidestyle Research Interests}
    % Computer vision, image processing, computational imaging, medical images, machine learning, embedded systems, the combination of computing and other areas.\\
 
     \vspace{0.2cm}
     \section{\mysidestyle Publication}
     \vspace{-0.35cm}
        \renewcommand{\refname}{}
        \renewcommand{\section}[1]{} % swallow the \section*{} emitted by thebibliography (it printed a stray "*")
        \bibliographystyle{ieeetr}
        \bibliography{publications}


     
        
     
\end{resume}   

 
\end{document}